% =====================================================================
%  GMN2 §5  Fock–Goncharov 坐标
%  原文：wkb-paper.tex 4062–4707 行；印刷页 p.63–75；编号公式 (5.1)–(5.40) 共 40 条
%  脚注 3 条；插图 9 幅（图 14–22，按原文 \insfig 出现次序数）
%  块内公式用 \eqref；跨块引用硬编码原文编号并加 % TODO 合稿改 \ref / \eqref
% =====================================================================

\section{Fock--Goncharov 坐标}
\label{sec:fg-coordinates}

接下来几节我们专门处理一个简单情形：设 $\CM$ 是 $C$ 上 Hitchin 方程解的模空间，规范群 $G=SU(2)$，在点 $\CP_1,\dots,\CP_{n_{\rm sing}}$ 处有 $n_{\rm sing}$ 个带半单留数的正则奇点。我们始终假定 $n_{\rm sing}\ge 1$；若 $C$ 的亏格为零，则假定 $n_{\rm sing}\ge 3$。\footnote{$n_{\rm sing}=3$、$C=\IC\IP^1$ 的情形有些退化，因为此时 $\CM$ 是零维的。} 对正则奇点的限制将在第 8 节% TODO 合稿改 \ref{sec:irregular}
中解除。

固定某个 $\zeta\in\IC^\times$。于是可以把 $\CM$（看作复结构 $J^{(\zeta)}$ 下的全纯辛流形）与平坦 $SL(2,\IC)$-联络 $\CA$ 的模空间等同起来。本节我们采用 Fock 与 Goncharov 的做法\cite{MR2233852}，在 $\CM$ 上定义一族有用的全纯达布坐标系。每个坐标系 $\CX^T$ 都联系于一个「装饰三角剖分」$T$——这是稍后就要定义的一种组合对象。随后几节我们将说明，如何由这些坐标系 $\CX^T$ 构造出 $\CM$ 上所要的函数 $\CX_\gamma$。

\subsection{Fock--Goncharov 坐标的定义}

所谓「三角剖分」，我们总是指 $C$ 的一个三角剖分，其顶点全部落在奇点 $\CP_i$ 处，且每个顶点至少关联一条边。在每个 $\CP_i$ 处有算子 $M_i$，给出 $\CA$-平坦截面绕 $\CP_i$ 的顺时针单值；它取值于 $SL(2,\IC)$，一般有两条互异的本征线。我们把 $\CP_i$ 处的一个\emph{装饰}定义为这两条本征线中选定一条；而一个\emph{装饰三角剖分} $T$ 就是一个三角剖分再加上每个顶点处的一个装饰。记 $\mu^T_i$ 为相应的单值本征值。

为避免混淆，值得指出我们的设置与\cite{MR2233852}有一处细微差别。他们的观点是把装饰的选取并入模空间，所以他们实际上是在「装饰平坦联络」的模空间上建立坐标系。而在我们的情形里，$M_i$ 的共轭类是一劳永逸地固定的，因此我们把装饰的选取算作离散数据 $T$ 的一部分。

\insfig{3in}{quadrilateral}{联系于三角剖分 $T$ 中边 $E$ 的四边形 $Q_E$。}

现在固定一个装饰三角剖分 $T$。对 $T$ 的每条边 $E$，如下定义一个坐标函数 $\CX^T_E$。界定 $E$ 的两个三角形拼成一个四边形 $Q_E$。按逆时针次序（用 $C$ 上的标准定向）把它的顶点编号为 $\CP_i$，$i=1,2,3,4$，并让 $E$ 连接顶点 $1$ 与 $3$。边是无向的，所以这样的标号只定到同时交换 $1\leftrightarrow 3$ 与 $2\leftrightarrow 4$；下面构造的 $\CX^T_E$ 与这一歧义无关。见图 14。% TODO 合稿改 \ref{fig:quadrilateral}

在 $Q_E$ 上我们选取 $V$ 的四个截面 $s_i$，使其满足平坦性方程
\begin{equation}
(d+\CA)s_i = 0,
\end{equation}
并且每个 $s_i$ 都是 $M_i$ 的本征向量，本征值为 $\mu^T_i$。这些 $s_i$ 不可能在 $C$ 上整体地单值且光滑（单值性会要求在某处引入割线），但我们强调，我们确实把它们取成在 $Q_E$ 上单值且光滑。每个 $s_i$ 只定到复标度因子。

Fock--Goncharov 坐标由这些 $s_i$ 构造：
\begin{equation}\label{eq:cross-ratio-wedges}
\CX^T_E := -\frac{(s_1\wedge s_2)(s_3\wedge s_4)}{(s_2\wedge s_3)(s_4\wedge s_1)},
\end{equation}
其中所有 $s_i$ 都在 $Q_E$ 中任取的同一个点 $\CP_*$ 处取值。由于联络 $\CA$ 取值于 $sl(2,\IC)$，这个量不依赖于 $\CP_*$。此外，每个 $s_i$ 的复标度歧义在 $\CX^T_E$ 中相互抵消。

函数 $\CX^T_E$ 只在 $s_i\wedge s_j=0$ 时（对两个相邻顶点 $\CP_i$、$\CP_j$）趋于 $0$ 或 $\infty$，而这发生在 $\CM$ 的一个余维 $1$ 子簇上。因此 $\CX^T_E$ 在 Zariski 开补片 $\CU_T\subset\CM$ 上有良好定义。事实上，$\CX^T_E$ 给出该补片上的一组坐标系\cite{MR2233852}：在附录 A% TODO 合稿改 \ref{app:MonodromyX}
中我们将说明，如何由 $\CX^T_E$ 重构出模掉规范等价的联络 $\CA$（即 $\CA$ 的单值表示）。

顺带一提，可以把 $\CX^T_E$ 看作交比，含义如下。在 $Q_E$ 上由 $\CA$-平坦截面构成的二维空间里取一组参照基之后，四个一维子空间 $\{\lambda s_i:\lambda\in\IC^\times\}\subset\IC^2$ 给出 $\IC\IP^1$ 上的四个点 $x_i$。于是可以证明，\eqref{eq:cross-ratio-wedges} 恰好就是
\begin{equation}\label{eq:cross-ratio}
\CX^T_E = -\frac{(x_1-x_2)(x_3-x_4)}{(x_2-x_3)(x_4-x_1)}.
\end{equation}

\subsection{单值}
\label{sec:monodromies}

$\CX^T_E$ 的某些组合有简单的解释。事实上，取任一奇点 $\CP_i$，考虑与 $\CP_i$ 相交的所有边 $E$ 上的 $\CX^T_E$ 之积。若不计 $s_i$ 的单值性，形式上会发现乘积中的各项望远镜式相消，只剩下 $1$。然而，定义 $\CX^T_E$ 时必须小心，让 $s_i$ 在整个四边形 $Q_E$ 上单值且连续。这样的选取可以对每条 $E$ 分别作出，但我们不可能构造出处处满足此条件的单一 $s_i$：总得在某处引入割线。若把割线取成穿过与 $\CP_i$ 相交的某一个三角形，则恰好有两个四边形受影响。它们各自乘上一个因子 $\mu^T$，于是得到
\begin{equation}\label{eq:monodromy-product}
\prod_{E\;\text{与}\;\CP_i\;\text{相交}} \CX^T_{E} = (\mu^T_i)^{2}.
\end{equation}
（我们暗中假定了 $\CP_i$ 至少关联两条边；实际上也可能只关联一条，此时 \eqref{eq:monodromy-product} 仍然成立，见下文第 5.8 节。）% TODO 合稿改 \ref{sec:degenerate}

\subsection{坐标的计数}

正如已经提到的，函数 $\CX^T_E$ 给出 $\CM$ 上的一组局部坐标系。作为一个简单的检验，我们来说明维数确实对得上。

$\CM$ 的维数可以由单值数据求得。设 $g$ 是 $C$ 的亏格，并回想我们的局部系统在 $C$ 上有 $n_{\rm sing}$ 个正则奇点 $\CP_i$。取值于 $SL(2,\IC)$ 的单值矩阵 $M_i$、$A_j$、$B_j$（分别绕奇点、$A$-圈与 $B$-圈）共有 $3(n_{\rm sing}+2g)$ 个自由度，减去 $3$（通常的约束 $\prod M_i=\prod A_jB_jA_j^{-1}B_j^{-1}$），再减去 $3$（$SL(2,\IC)$ 规范对称），再减去 $n_{\rm sing}$（因为单值本征值固定）。合起来给出
\begin{equation}
\dim\CM = 2n_{\rm sing}+6g-6.
\end{equation}
我们又能得到多少个坐标 $\CX^T_E$？利用 Euler 公式
\begin{equation}
n_{\rm sing}+(\#F)-(\#E)=2-2g
\end{equation}
以及所有面都是三角形这一事实
\begin{equation}
2(\#E)=3(\#F)
\end{equation}
得到
\begin{equation}\label{eq:counts}
\#E=3n_{\rm sing}+6g-6,\qquad \#F=2n_{\rm sing}+4g-4.
\end{equation}
如上面所论证，$\CX^T_E$ 中有 $n_{\rm sing}$ 个组合给出单值本征值 $(\mu^T_i)^2$，而这些本征值只定到一种离散的选择；其余的 $2n_{\rm sing}+6g-6$ 个恰好足以给出 $\CM$ 上的坐标，正如所愿。

\subsection{哈密顿流}
\label{sec:ham-flows}

正如我们在第 4 节% TODO 合稿改 \ref{sec:hitchin}
中所指出的，$\CM$ 带有一个自然的全纯辛形式。现在我们来描述由 $\log \CX^T_E$ 生成的哈密顿流。这个流原来非常简单。

给定一个联络 $\CA\in\CM$，可以如下描述它在 $\log\CX^T_E$ 生成的时间-$t$ 流之下的像 $\CA^t$。把 $C$ 上的联络 $\CA$ 看成两块，即限制到四边形 $Q_E$ 上的部分（「内部」）与限制到其补集 $C\setminus Q_E$ 上的部分（「外部」）。在两块的公共边界上，我们有一个 $V$ 的截面之边界值之间的等同。用这个等同把两块粘回去，就还原出原来的 $V$ 及其联络 $\CA$。现在考虑换一个等同来粘。具体说，在边 $E_{ij}$ 上（$i$、$j$ 按循环次序），我们把外部的 $s_i$ 粘到内部的 $e^{\pm t/2}s_i$，把外部的 $s_j$ 粘到内部的 $e^{\mp t/2}s_j$，其中符号 $\pm$ 当 $i=2,4$ 时取 $+$，当 $i=1,3$ 时取 $-$。这个粘合仍取值于 $SL(2,\IC)$，因为它保持 $s_i\wedge s_j$；它也保持联络的平坦性，因为我们把平坦截面粘到了平坦截面上。它定义了新联络 $\CA^t$。

注意截面 $s_i$ 绕 $\CP_i$ 的单值对 $\CA^t$ 与对 $\CA$ 是一样的（因为来自与 $\CP_i$ 相交的两条边的 $s_i$ 变换相互抵消）。由此可知 $\CA^t$ 与 $\CA$ 有相同的单值本征值，所以这个流确实保持我们的模空间 $\CM$，理当如此。

由辛结构 (4.10)% TODO 跨块引用
推出这个流基本上是直截了当的，只是记号上有些笨拙，所以我们把它挪到了附录 B% TODO 合稿改 \ref{app:poorman}
中。

\subsection{泊松括号}
\label{sec:poisson}

给定三角剖分 $T$ 的两条边 $E$ 与 $E'$，定义 $\inprod{E,E'}$ 为 $E$ 与 $E'$ 公共面的个数，计数时带符号：若绕公共面按逆时针（顺时针）方向走时 $E$ 紧接在 $E'$ 之前，则取 $+1$（$-1$）。本节我们证明，$\CX^T_E$ 的泊松括号由这个配对所决定，
\begin{equation}\label{eq:fg-poisson}
\{\CX^T_E,\CX^T_{E'}\} = \inprod{E,E'}\,\CX^T_E\,\CX^T_{E'}.
\end{equation}

为了验证 \eqref{eq:fg-poisson}，我们考虑由 $\log\CX^T_E$ 生成的哈密顿流对 $\CX^T_{E'}$ 的作用。回想第 5.4 节，% TODO 合稿改 \ref{sec:ham-flows}
这个流是沿四边形 $Q_E$ 的各边作切割与粘合。若 $\inprod{E,E'}=0$，则只用 $Q_E$ 外面的联络 $\CA$ 就能算出 $\CX^T_{E'}$，因此 $\CX^T_{E'}$ 在该流下不变，与 \eqref{eq:fg-poisson} 相符。现在设 $\inprod{E,E'}=+1$。定义 \eqref{eq:cross-ratio-wedges} 要求我们把截面 $s_i$ 从 $Q_{E'}$ 的顶点 $\CP_{1,2,3,4}$ 平移到一个公共点 $\CP_*\in Q_{E'}$；让我们把这个点也取在 $Q_E$ 内。顶点标号的约定见图 15。% TODO 合稿改 \ref{fig:poisson-bracket}

\insfig{3.5in}{poisson-bracket}{为计算由 $\log\CX^T_E$ 生成的流对函数 $\CX^T_{E'}$ 的作用而固定记号，情形为 $\inprod{E,E'}=+1$。}

这一平移要求我们把 $s_5$ 从 $Q_E$ 外带到内部，跨过边 $E'=E_{12}$。由此可见，定义 $\CX^T_{E'}$ 时出现的 $s_5$ 会被该流变换。把 $s_5$ 在方便的平坦截面基下展开为
\begin{equation}
s_5 = a s_1 + b s_2
\end{equation}
则流的作用是 $a\mapsto ae^{-t/2}$、$b\mapsto be^{t/2}$，从而
\begin{equation}
s_1\wedge s_5 \to e^{t/2} s_1\wedge s_2,
\end{equation}
\begin{equation}
s_5\wedge s_2 \to e^{-t/2} s_5\wedge s_2.
\end{equation}
于是，由 $\log\CX^T_E$ 生成的流把 $\CX^T_{E'}$ 乘以 $e^{t}$。所要的泊松括号 \eqref{eq:fg-poisson} 直接得到。

也可能出现 $\inprod{E,E'}=+2$，即两条边共享两个面。此时类似的论证表明，由 $\log\CX^T_E$ 生成的流把 $\CX^T_{E'}$ 乘以 $e^{2t}$，所要的 \eqref{eq:fg-poisson} 仍然成立。

\subsection{坐标变换与装饰三角剖分的群胚}

现在到了关键之处。坐标系 $\CX^T$ 依赖于三角剖分 $T$ 的选取。我们想知道把 $\CX^T$ 与 $\CX^{T'}$ 联系起来的坐标变换 $S_{T,T'}$。

有一个计算 $S_{T,T'}$ 的程序，原则上直截了当，实践中却相当繁琐。交比 $\CX^T_E$ 含有足够的信息，可以沿 $C$ 中的路径平移平坦截面 $s$，见附录 A。% TODO 合稿改 \ref{app:MonodromyX}
因此，若 $T$ 与 $T'$ 有相同的装饰，就可以先用 $\CX^T_E$ 定出沿 $T'$ 的边 $E'_{ij}$ 的 $s_i$ 与 $s_j$，再用它们算出 $\CX^{T'}_{E'_{ij}}$。若装饰不同，则还需要确定新的单值本征截面 $s'_i$，办法是把参照截面绕三角剖分 $T$ 的顶点作平移。这样得到的 $S_{T,T'}$ 的显式表达式一般是复杂的有理变换，对我们用处不大（只有一个关键例外，我们将在第 11 节% TODO 合稿改 \ref{sec:UniversalStokesMatrix}
遇到）。

与其如此，更有用的是研究对应于三角剖分 $T$ 的简单局部改变的坐标变换。任意两个三角剖分 $T$、$T'$ 都可以通过一列这样的简单变换互相联系，我们称之为「边上的 \emph{flip}」（翻转一条边）与「顶点处的 \emph{pop}」（改换该顶点处本征线的选取）：

\begin{itemize}
\item 给定 $T$ 的一条边 $E$，如下定义 flip $\sigma_E$：原来的 $T$ 与 flip 后的 $T'$ 完全相同，只是在四边形 $Q_E$ 中，$T$ 的边 $E=E_{13}$ 被换成 $T'$ 的边 $E'=E_{24}$，如图 16% TODO 合稿改 \ref{fig:flip}
所示。（这里假定 $E$ 不同于四边形 $Q_E$ 的所有边。当存在两条边重合的退化三角形时，这一点可能不成立，下面会讨论。对这样的边我们不需要定义 flip。）
\insfig{5.5in}{flip}{把三角剖分 $T$ flip 成 $T'$。}

\item 给定顶点 $\CP_i$，如下定义 pop $\pi_i$：原来的 $T$ 与 pop 后的 $T'$ 只相差在 $\CP_i$ 处反转单值本征值（装饰）的选取。
\end{itemize}

于是，坐标变换 $S_{T,T'}$ 可以方便地描述为一列把 $T$ 带到 $T'$ 的 flip 与 pop 所对应的坐标变换之复合。

用群胚的语言重新表述上述内容是方便的。\footnote{「群胚」是一个所有态射都可逆的范畴。用白话讲，这只是说我们有一套点（「对象」）以及点之间的一族箭头（「态射」）：箭头上有一个结合的复合律，每个点上有一个单位箭头，每个箭头都有一个逆箭头。} 首先设想一个群胚 $\CT$，其对象是各种装饰三角剖分，而任意两个装饰三角剖分之间恰有一个态射。（也就是说，对象没有自同构。）原来 $\CT$ 由 flip 与 pop 自由生成，只受四类关系的约束：
\begin{itemize}
\item 当 $\inprod{E,E'}=0$ 时，$\sigma_E\sigma_{E'}=\sigma_{E'}\sigma_E$。

\item 当 $\inprod{E,E'}=1$ 时，$\sigma_E\sigma_{E'}\neq\sigma_{E'}\sigma_E$，但 $\sigma_E\sigma_{E'}$ 可以改写成\emph{三个} flip 之积，如图 17% TODO 合稿改 \ref{fig:pentagon-flips}
所示。\footnote{当 $\inprod{E,E'}=2$ 时，$\sigma_E$ 与 $\sigma_{E'}$ 之间没有关系，但这一情形在下文第 5.9 节% TODO 合稿改 \ref{sec:limit-triangulation}
中仍将很重要。}

\item 每个 $\pi_i$ 与其余一切都对易。

\item $\pi_i^2=\mathrm{id}$。

\end{itemize}

\insfig{4in}{pentagon-flips}{「五边形关系」：联系一对三角剖分的两种不同的 flip 序列。}

其次，定义 $\CM$ 上达布坐标系的群胚 $\CC$。为此先定义一个\emph{泊松环面}，即同构于 $(\IC^\times)^{\dim\CM}$ 的空间，带一个泊松结构，使得标准坐标函数 $x_i$ 的括号形如 $\{x_i,x_j\}=a_{ij}x_ix_j$。于是一个\emph{达布坐标系}就是从 $\CM$ 的某个 Zariski 开子集到某个泊松环面的单射泊松映射。群胚 $\CC$ 以达布坐标系为对象，任意两个对象之间恰有一个态射，可等同于它们之间唯一的达布坐标变换（即重叠区在 $\CM$ 中的两个像之间的一个泊松双射）。

于是可以把 \eqref{eq:cross-ratio-wedges} 看成在定义一个函子，
\begin{equation}
\CT \to \CC
\end{equation}
\begin{equation}
T \mapsto \CX^T.
\end{equation}
在这个函子下，三角剖分的态射被映为相应的坐标变换，所以态射之间的关系蕴涵坐标变换之间的关系。我们马上就会看到，这些关系写成显式时可能出人意料。

\subsection{flip 下的变换}
\label{subsec:TransfmFlips}

一个 flip $\sigma_E$ 对坐标的作用相当简单。首先，参照图 16% TODO 合稿改 \ref{fig:flip}
容易看出，$T$ 中附着于 $E$ 的坐标与 $T'$ 中附着于 $E'$ 的坐标之间有平凡的关系，
\begin{equation}\label{eq:fg-trans-1}
\CX^T_E = (\CX^{T'}_{E'})^{-1}.
\end{equation}

\insfig{5in}{flip-edge-coordinate}{flip 对用于计算 $\CX^T_{E_{41}}$ 与 $\CX^{T'}_{E_{41}}$ 的四边形的影响。}

flip 也改变附着于 $Q_E$ 四条边的坐标。例如，考虑图 18% TODO 合稿改 \ref{fig:flip-edge-coordinate}
中的坐标 $\CX^T_{E_{41}}$。它是 $s_1$、$s_3$、$s_4$、$s_5$ 依此次序的交比，我们记作 $\CX^T_{E_{41}}=r(1,3,4,5)$。作 flip 之后，代之以 $\CX^{T'}_{E_{41}}=r(1,2,4,5)$。另一方面，我们还有 $\CX^T_E=r(1,2,3,4)$。既然 $\IC\IP^1$ 上任意五个点只有两个独立的交比，$\CX^T_{E_{41}}$、$\CX^{T'}_{E_{41}}$ 与 $\CX^T_E$ 之间就必定存在一个代数关系。直接计算确实给出这个关系，以及涉及 $Q_E$ 其余边的类似关系：
\begin{equation}\label{eq:fg-trans-2}
\CX^{T'}_{E_{12}} = \CX^T_{E_{12}}(1+\CX^T_E),
\end{equation}
\begin{equation}\label{eq:fg-trans-3}
\CX^{T'}_{E_{23}} = \CX^T_{E_{23}}(1+(\CX^T_E)^{-1})^{-1},
\end{equation}
\begin{equation}\label{eq:fg-trans-4}
\CX^{T'}_{E_{34}} = \CX^T_{E_{34}}(1+\CX^T_E),
\end{equation}
\begin{equation}\label{eq:fg-trans-5}
\CX^{T'}_{E_{41}} = \CX^T_{E_{41}}(1+(\CX^T_E)^{-1})^{-1}.
\end{equation}

（上面我们假定了每个 $\inprod{E_{ij},E}=\pm1$，即唯一的公共面就是图中出现的那个。若例如 $\inprod{E,E_{41}}=2$，则变换改为 $\CX^{T'}_{E_{41}}=\CX^T_{E_{41}}(1+(\CX^T_E)^{-1})^{-2}$，其余边类似。）

这些变换当然保持泊松括号 \eqref{eq:fg-poisson}（保持的方式并不平凡，因为相交矩阵 $\inprod{E,E'}$ 在 flip 下会变）。所以它们定义了达布坐标系的一个态射。

此外要注意，这个态射与我们在第 2 节% TODO 合稿改 \ref{sec:review}
回顾过的 KS 变换 $\CK_\gamma$ 十分相似。这是 Fock--Goncharov 坐标与 Kontsevich--Soibelman 穿墙公式之间联系的第一条线索。不过两者的变换律还不完全相同；要作出认同，还需要下面就要描述的重要一步。

最后我们指出：装饰三角剖分的群胚 $\CT$ 中的五边形恒等式，蕴涵 $\CC$ 中态射之间的一个相应五边形恒等式。它具体是怎样产生的？再来看图 17，% TODO 合稿改 \ref{fig:pentagon-flips}
把第 $n$ 个五边形两条内边所对应的两个坐标函数记作 $(x_n,y_n)$，其中 $n$ 按模 $5$ 计。联系相邻两个五边形的 flip 给出坐标变换
\begin{equation}
y_{n+1} = x_n^{-1},
\end{equation}
\begin{equation}
x_{n+1} = y_n(1+x_n).
\end{equation}
这一列坐标变换确实以 $5$ 为周期。特别地，消去 $y_n$ 就得到漂亮而简单的 5 周期关系 $x_{n+1}x_{n-1}=1+x_n$。

\subsection{退化三角剖分与 pop 下的变换}
\label{sec:degenerate}

对一般三角剖分 $T$ 的一个顶点 $\CP$ 作 pop，效果通常相当复杂。要用旧的平坦截面构造出 $\CP$ 处新的平坦截面，我们需要绕 $\CP$ 整整一圈作平行移动。原则上这由关联到 $\CP$ 的那些边上的 $\CX^T_E$ 决定，但实践中结果一般是复杂的有理函数。本小节我们将遇到一类特殊的三角剖分：只有\emph{一条}边关联到 $\CP$，此时 pop 的效果变得非常简单。这些三角剖分含有退化的面，其中两条边被等同，如图 19% TODO 合稿改 \ref{fig:degenerate-face}
所示。

当含有退化面时，我们构造坐标 $\CX^T_E$ 的规则须稍作修正。我们通过过渡到某个邻域 $U\subset C$ 的覆叠 $\tilde U$ 来「化解」这个面，该覆叠只在中心顶点处分歧。$U$ 上的平坦联络 $\CA$ 拉回到 $\tilde U$ 上的平坦联络，在 $U$ 上各奇点的原像处有正则奇点，而 $T|_U$ 提升为三角剖分 $\tilde T$。我们选取覆叠使 $\tilde T$ 不退化，于是可以用标准规则定义坐标 $\CX^{\tilde T}_{\tilde E}$。（为了完全化解这个面，覆叠至少应有三层，如图 20% TODO 合稿改 \ref{fig:degenerate-quadrilateral}
所示。）然后定义 $\CX^T_E:=\CX^{\tilde T}_{\tilde E}$，其中 $\tilde E$ 是 $E$ 的任一原像。由于联络是拉回的，它对覆叠的自同构不变，所以 $\CX^T_E$ 与 $\tilde E$ 的选取无关；同理，它与覆叠的确切选取也无关。

\insfig{1.25in}{degenerate-face}{一个退化面：可以把它想成一个三角形，其三边为 $E'$、$E$ 与 $E$。}

\insfig{4.5in}{degenerate-quadrilateral}{退化面的一个邻域，以及通过过渡到三重覆叠所作的化解。}

例如，我们利用四边形 $Q_{\tilde E_1}$ 如下算出 $\CX^T_E$。设 $s$ 是 $\CP$ 的原像 $\tilde\CP$ 处的装饰，我们可以把它取成 $\CP$ 处某个装饰的拉回。设 $s'_i$ 是 $\CP'$ 的各原像 $\tilde\CP'_i$ 处的装饰；我们可以把所有 $s'_i$ 都取成 $\CP'$ 处同一个装饰的拉回。于是 \eqref{eq:cross-ratio-wedges} 变成
\begin{equation}
\CX^T_E = \CX^{\tilde T}_{\tilde E_1} = -\frac{(s'_3\wedge s'_1)(s'_2\wedge s)}{(s\wedge s'_3)(s'_1\wedge s'_2)},
\end{equation}
它可以进一步化简如下。设 $M$ 表示原退化面中绕 $\CP$ 的顺时针单值。把所有 $s'_i$ 平移到 $Q_{\tilde E_1}$ 中的公共点后，我们将有 $s'_1=M^{-1}s'_3=Ms'_2$，于是
\begin{equation}\label{eq:degenerate-monodromy}
\begin{split}
\CX^T_E &= -\frac{(Ms'_1\wedge s'_1)(M^{-1}s'_1\wedge s)}{(s\wedge Ms'_1)(s'_1\wedge M^{-1}s'_1)}\\
&= -\frac{(Ms'_1\wedge s'_1)(s'_1\wedge Ms)}{(M^{-1}s\wedge s'_1)(Ms'_1\wedge s'_1)} = (\mu^T)^{2},
\end{split}
\end{equation}
其中回想 $\mu^T$ 是 $s$ 在 $M$ 下的本征值，并且我们用到了 $M\in SL(2,\IC)$。

对 $\CX^T_{E'}$，我们得到
\begin{equation}
\CX^T_{E'} = \CX^{\tilde T}_{\tilde E'_3} = -\frac{(s\wedge s'_3)(s''_3\wedge s'_1)}{(s'_1\wedge s)(s'_3\wedge s''_3)},
\end{equation}
同样，在平移到 $Q_{\tilde E'_3}$ 中的公共点后有 $s'_3=Ms'_1$，于是
\begin{equation}\label{eq:degen-loop-coord}
\begin{split}
\CX^T_{E'} &= -\frac{(s\wedge Ms'_1)(s''_3\wedge s'_1)}{(s'_1\wedge s)(Ms'_1\wedge s''_3)}\\
&= -\frac{M^{-1}s\wedge s'_1}{s'_1\wedge s}\,\frac{s''_3\wedge s'_1}{Ms'_1\wedge s''_3}
= (\mu^T)^{-1}\,\frac{Ms'_1\wedge s''_3}{s''_3\wedge s'_1}.
\end{split}
\end{equation}
注意 $s$ 只通过其单值本征值 $\mu^T$ 进入这个结果。由此可见，在顶点 $\CP$ 处作 pop 的唯一效果，就是通过关系 $\mu^{T'}=(\mu^T)^{-1}$ 起作用；利用 \eqref{eq:degenerate-monodromy}、\eqref{eq:degen-loop-coord}，这给出变换 $\pi_\CP$：
\begin{equation}
\CX^{T'}_E = (\CX^{T}_E)^{-1},
\end{equation}
\begin{equation}
\CX^{T'}_{E'} = \CX^T_E\,\CX^T_{E'}.
\end{equation}

所以，在退化顶点 $\CP$ 处作 pop 的效果非常简单。另一方面，对一般三角剖分 $T$ 在 $\CP$ 处作 pop 的效果也可确定：先把关联到 $\CP$ 的边中除一条之外全部 flip，达到只有一条边关联到 $\CP$ 的退化三角剖分，再在 $\CP$ 处作 pop，最后按相反次序把各边 flip 回去。由于 pop 与 flip 对易，这与直接在 $T$ 的 $\CP$ 处作 pop 相同。更一般地，对一般的一对三角剖分 $T$、$T'$ 计算 $S_{T,T'}$ 的一个办法，是把该态射分解为一列 flip：这些 flip 把 $T$ 底下的无装饰三角剖分变成 $T'$ 底下的那个，并途经退化三角剖分——在退化三角剖分处，所需 pop 的效果是简单的。

\subsection{三角剖分的极限与 \emph{juggle}}
\label{sec:limit-triangulation}

到此为止，我们考虑了坐标 $\CX^T_E$ 在 flip 与 pop 下的变换。为我们的目的，实际上还必须考虑第三种更复杂的变换：它并不完全是两个三角剖分之间的关系，而是三角剖分的两种「无限扭曲极限」之间的关系。本小节我们引入这些极限。

我们将考虑含有环形区域 $W$ 的三角剖分 $T$，外圈上有一个顶点 $\CP$，内圈上有一个顶点 $\CP'$。任何这样的 $T$ 在 $W$ 上有两条内部边，且两条边都以这两个顶点为端点。假定我们固定 $T$ 在 $W$ 之外的部分，而考虑改变 $W$ 上的部分。这类 $T$ 有多种选择，彼此的区别在于边绕环缠了多少圈。图 21% TODO 合稿改 \ref{fig:annulus}
给出几个例子。

为了参数化可能的 $T$，我们先选定两条固定路径 $E_\pm$，从 $\CP$ 到 $\CP'$，使得 $E_+-E_-$ 绕 $W$ 逆时针转一圈。然后定义 $T_0$ 为内部边是 $E_{0+}:=E_+$、$E_{0-}:=E_-$ 的三角剖分。

对 $E_{0-}$ 作一次 flip，我们得到另一个三角剖分 $T_1$。把它的边标为 $E_{1\pm}$，同样让 $E_{1+}$ 与 $E_{1-}$ 相差一个单位的逆时针缠绕：于是 $E_{1+}$ 是 flip 新造出的那条边，而 $E_{1-}=E_{0+}$。再对 $E_{1-}$ 作 flip，又得到新三角剖分 $T_2$。见图 21。% TODO 合稿改 \ref{fig:annulus}
重复这一过程，就得到对 $m\in\IZ_+$ 的一列三角剖分 $T_m$。flip $E_{m-}$ 把 $T_m$ 带到 $T_{m+1}$；反之，flip $E_{m+}$ 把 $T_m$ 带到 $T_{m-1}$。

\insfig{13cm}{annulus}{三角剖分 $T_0$ 中的一个环，以及由 flip $E_{0-}$、$E_{1-}$ 得到的三角剖分 $T_1$、$T_2$。}

三角剖分 $T_m$ 的 Fock--Goncharov 坐标是什么？为了尽量减少混淆，我们过渡到环的某个无限覆叠，就像讨论退化三角剖分时所用的那些覆叠一样，并选定 $E_\pm$ 的特定原像 $\tilde E_\pm$。定义 $s$、$s'$ 为 $\tilde E_+$ 两端 $\tilde\CP$、$\tilde\CP'$ 处的装饰，并引入记号
\begin{equation}
K = -(s\wedge Ms)(s'\wedge Ms'),\qquad c_k = (s\wedge M^k s')^2.
\end{equation}
定义 \eqref{eq:cross-ratio-wedges} 于是变成（见图 22）% TODO 合稿改 \ref{fig:log-cover-Tm}
\begin{equation}
\CX^{T_m}_{E_{m+}} = \frac{K}{c_{1-m}},\qquad \CX^{T_m}_{E_{m-}} = \frac{c_{-m}}{K}.
\end{equation}

\insfig{11cm}{log-cover-Tm}{我们对环所作的无限覆叠，并标出三角剖分 $T_m$。我们定义 $\tilde\CP_0:=\tilde\CP$、$\tilde\CP'_0:=\tilde\CP'$，然后把它们的平移分别记作 $\tilde\CP_m$ 与 $\tilde\CP'_m$。}

我们这里主要关心 $m\to\infty$ 的极限。如下定义两个「极限坐标」$\CX^{T_{+\infty}}_{A,B}$。记 $M$ 的本征值为 $\xi_\pm$，其中 $|\xi_+|>1$。（我们假定处于一般情形，即 $|\xi_\pm|\neq1$。在物理应用中，这相当于假定矢量多重态的质量是一般的。）其中一个坐标就是
\begin{equation}\label{eq:A-positive}
\CX^{T_{+\infty}}_A := \xi_+^{2}.
\end{equation}
为定义另一个坐标，引入两条单值本征空间上的投影算子 $P_\pm$。然后定义
\begin{equation}\label{eq:B-positive}
\CX^{T_{+\infty}}_B := -\frac{(s\wedge P_- s')^2}{(s\wedge Ms)(s'\wedge Ms')}.
\end{equation}
直接计算表明，它们是「极限」坐标，意思是
\begin{equation}\label{eq:limiting-1}
\CX^{T_{+\infty}}_A = \lim_{m\to\infty}\CX^{T_m}_{E_{m+}}\CX^{T_m}_{E_{m-}},
\end{equation}
\begin{equation}\label{eq:limiting-2}
\CX^{T_{+\infty}}_B = \lim_{m\to\infty}(\CX^{T_m}_{E_{m+}})^{-m}(\CX^{T_m}_{E_{m-}})^{1-m}.
\end{equation}
令 $E$ 表示环的某一条边界，则坐标 $\CX^{T_m}_E$ 有良好定义的 $m\to\infty$ 极限，我们称之为 $\CX^{T_{+\infty}}_E$。环外所有边的坐标根本不依赖于 $m$，所以对这类边 $E$，我们直接定义 $\CX^{T_{+\infty}}_E=\CX^{T_m}_E$（对任意 $m$）。于是可以定义新的坐标系 $\CX^{T_{+\infty}}$，它由 $\CX^{T_{+\infty}}_A$、$\CX^{T_{+\infty}}_B$ 以及所有其余的 $\CX^{T_{+\infty}}_E$ 组成。

我们在这里得到的坐标系 $\CX^{T_{+\infty}}$ 依赖于一开始所作的选取，即把哪两条路径叫做 $E_\pm$。$\CX^{T_{+\infty}}_A$ 与这一选取无关，但不同选取下的 $\CX^{T_{+\infty}}_B$ 相差 $\xi_+$ 的整数次幂。

我们也可以考虑相反的一类极限。再次从 $T$ 出发（并选定同样的路径 $E_\pm$），反复 flip $E_{m+}$，就得到 $m<0$ 的一列三角剖分 $T_m$，其边沿逆时针方向绕环。然后类似地定义
\begin{equation}\label{eq:A-negative}
\CX^{T_{-\infty}}_A := \xi_-^2
\end{equation}
以及
\begin{equation}\label{eq:B-negative}
\CX^{T_{-\infty}}_B := -\frac{(s\wedge P_+ s')^2}{(s\wedge Ms)(s'\wedge Ms')}.
\end{equation}
它们也是极限坐标，意思是
\begin{equation}\label{eq:limiting-3}
\CX^{T_{-\infty}}_A = \lim_{m\to-\infty}\CX^{T_m}_{E_{m+}}\CX^{T_m}_{E_{m-}},
\end{equation}
\begin{equation}\label{eq:limiting-4}
\CX^{T_{-\infty}}_B = \lim_{m\to-\infty}(\CX^{T_m}_{E_{m+}})^{-m}(\CX^{T_m}_{E_{m-}})^{1-m}.
\end{equation}

从群胚 $\CT$ 的观点看，我们是在引入可数多个新对象：对参照路径 $E_\pm$ 的可数多种可能选取，各加上两个「极限三角剖分」$T_{\pm\infty}$。相对于某个固定的 $E_\pm$ 选取，可以把这些极限三角剖分更具体地标为 $T_{\pm\infty}^{[k]}$，$k\in\IZ$。存在从 $T_m$ 到任一 $T_{\pm\infty}^{[k]}$ 的态射，它表示一列无限多个 flip 的极限。还有一个相对平凡的、改变 $k$ 的态射：在坐标系的层面上，它让 $\CX_A$ 不变，而把 $\CX_B$ 改变 $(\xi_+)^{k'-k}$ 倍。我们还引入更有意思的态射，称之为「\emph{juggle}」（在两类极限三角剖分之间跳接的态射）：从任一 $T^{[k]}_{+\infty}$ 到任一 $T^{[k']}_{-\infty}$ 有一个 juggle 态射，反向亦然。从道德上说，引入这些 juggle 是在三角剖分的群胚中「绕出一个闭环」：从任一 $T_m$ 出发，我们可以想象 flip 无限多次到达某个 $T_{+\infty}$，再 juggle 到某个 $T_{-\infty}$，然后再 flip 无限多次回到 $T_m$。我们将滥用记号，仍用 $\CT$ 表示由这些极限三角剖分与额外态射增广后的装饰三角剖分群胚。

那么 juggle 如何作用于坐标系，即它在我们函子 $\CT\to\CC$ 下的像是什么？对 $A$ 坐标这很简单：由 \eqref{eq:A-positive} 与 \eqref{eq:A-negative} 立即得到
\begin{equation}\label{eq:A-relation}
\CX^{T_{-\infty}}_A = (\CX^{T_{\infty}}_A)^{-1}.% 原文如此：右端指数印作 T_{\infty}（据上下文应为 T_{+\infty}），照译
\end{equation}
$B$ 坐标又如何呢？这里的关系取决于我们在定义 $T_{+\infty}$ 与 $T_{-\infty}$ 时对路径 $E_\pm$ 的选取。假定两者选用同一组 $E_\pm$。那么一个简单的线性代数计算表明
\begin{equation}
\frac{(s\wedge Ms)(s'\wedge Ms')}{(s\wedge P_+ s')(s\wedge P_- s')} = -(\xi_+-\xi_-)^2.
\end{equation}
把它与定义 \eqref{eq:B-positive}、\eqref{eq:B-negative} 结合，得到
\begin{equation}\label{eq:B-relation}
\CX^{T_{-\infty}}_B = (\CX^{T_{+\infty}}_B)^{-1}(\xi_+-\xi_-)^{-4}.
\end{equation}
